\chapter{关于认真贯彻落实《保障中小企业款项支付条例》进一步做深做实清理拖欠中小企业账款工作的通知}

\begin{center}
国资发财评〔2021〕104号
\end{center}

\noindent 各中央企业：

当前我国经济发展面临需求收缩、供给冲击、预期转弱三重压力，中小企业在经营成本、市场需求、要素保障、政策环境、资金回笼等方面面临的困难和挑战明显增多，已成为重大宏观经济问题。党中央、国务院对此高度重视，中央经济工作会议、国务院常务会议作出重要部署，要求严格执行《保障中小企业款项支付条例》（以下简称《条例》），强化契约精神，有效治理恶意拖欠账款。为贯彻落实党中央、国务院决策部署，帮助纾解中小企业困难，现就进一步做深做实清理拖欠中小企业账款（以下简称清欠）工作有关事项通知如下：

\section{进一步提高政治站位，切实增强做好清欠工作的积极性主动性}

各中央企业要认真学习领会习近平总书记关于支持中小企业发展的系列重要讲话精神，增强“四个意识”，坚定“四个自信”，做到“两个维护”，立足新发展阶段，完整、准确、全面贯彻新发展理念，自觉将清欠工作纳入服务构建新发展格局、实现高质量发展工作中统筹考虑、系统谋划，以更加坚定的思想自觉、精准务实的举措、真抓实干的劲头，切实做好及时支付中小企业账款、支持中小企业发展工作，积极构建大中小企业相互依存、相互促进的企业发展生态，确保党中央、国务院决策部署落到实处，为做好“六稳”“六保”工作发挥积极作用。

\section{带头贯彻落实《条例》，加快健全完善防范拖欠长效机制}

各中央企业要将《条例》作为普法学习的重要内容，持续深入开展宣贯学习，强化契约精神，加强社会责任建设，在做好清欠工作基础上积极开展正面宣传。对照《条例》，按照《关于进一步巩固清欠工作成果 加快健全防止拖欠长效机制的通知》（国资发财评〔2020〕43号）等要求，认真检查本企业贯彻落实方案（计划）实施情况，全面梳理合同管理、支付流程、信息化管理等工作中与《条例》不相符的地方，强化工作保障，确保防范拖欠长效机制不断健全完善，从制度、机制、流程和信息化管控上杜绝滥用市场优势地位恶意拖欠账款行为。同时，要做好清欠工作与投资并购、资产重组、亏损治理等工作的有机统筹，对并购重组等带入的民营企业中小企业无分歧欠款，要纳入前期尽调和后期工作方案，并限期清偿。

\section{采取有力有效举措，进一步加强应付账款和应付票据管控}

各中央企业要针对当前应付账款和应付票据管控存在的问题，推进应收、应付“一起管”，努力实现集团整体应付账款和应付票据增幅低于营业成本增幅。一是坚持按合同办事，严格按照合同约定的时间、方式及时足额支付款项，同时积极运用信息化手段规范应付账款特别是中小企业账款支付。二是加强合规管理，清理霸王条款，不得滥用市场优势地位设立不合理的付款条件、时限，占压中小企业资金。三是严控“背靠背”付款条款，对于提前明示、合同约定“背靠背”付款条款的，要加强上游款项催收，上游付款后及时对中小企业付款。\cl{四是严格票据等非现金支付管理，新签合同要明确约定支付方式和时限，不得利用优势地位强迫中小企业接受非现金支付；}新签合同未事先书面约定非现金支付的，事后原则上不得通过补充协议等方式约定或使用非现金支付，确需事后使用非现金支付的，应当承担资金成本；2022年起，除原有合同已有书面约定外，\cl{原则上不再开具6个月以上的商业承兑汇票和供应链债务凭证，防止变相延长付款时限。}

\section{规范供应链金融业务，助力缓解中小企业“融资难”“融资贵”}

\cl{各中央企业要积极发挥产业链“核心”企业作用，支持配合上下游中小企业开展供应链融资，努力实现自身优质信用与上下游中小企业共享，助力缓解中小企业“融资难”“融资贵”。一是正确认识供应链金融作用，在依法合规、能力可及、风险可控的基础上，积极稳妥开展供应链金融业务，全面清理与国家有关支持民营企业中小企业精神不符的内部规章制度。二是坚持开放共享，中小企业以其持有的集团内单位出具的商票、供应链债务凭证、应付账款办理贴现、保理等融资业务的，不得拖延确权，不得将融资机构限定为集团内部机构（平台，下同），严禁高息套利。三是集团对相关内部机构绩效考核时，应降低上述业务利润权重或予以剔除，引导相关内部机构提升对上下游中小企业的服务质量，降低服务成本。}

\section{扎实做好线索核实，及时有效回应中小企业诉求}

各中央企业要将拖欠线索核实处理工作，作为纾解中小企业困难、检验《条例》落实成效的重要抓手，慎终如始持续做好拖欠线索核实工作。一是层层落实责任。集团和各级子企业要将清欠投诉举报方式公开在官网或公众号等醒目位置；在业务发生前或在合同中书面告知本企业及上级企业拖欠线索受理方式，引导中小企业优先向本企业及上级企业反映问题；认真核实属地政府部门移交的投诉，不能以“不归地方管”为由推诿塞责，尽可能把矛盾和纠纷解决在业务发生地、解决在基层一线，减少多头、层层、重复举报。二是坚持依法合规。核实处理线索要以法律规定和合同约定为依据，该付的必须及时付、坚决付，同时避免误导中小企业将投诉举报作为催款手段；有分歧的积极协商解决，不得以有分歧为由置之不理；已有生效司法、仲裁结果的，要及时支付；中小企业确有困难，需要提前付款的，要严格履行内部审批程序。三是提升质量效率。对各方面反映问题线索，要限期核实处理，处理结果及时反馈，不得简单以对方联系不上、未提交付款资料等为由敷衍了事；确属拖欠或存在分歧的，不得强行要求对方出具不拖欠“澄清函”“说明函”，有分歧事项应有外部或书面证据。四是严格审核把关。要加强对子企业线索核实质量和效率的监督检查，对国办“互联网+”督查、违约拖欠中小企业款项登记（投诉）平台等转来的线索，要逐级复核，层层压实责任，不能只当“二传手”。五是完善工作闭环。线索核实中发现拖欠无分歧账款、合同约定不合理不明确、投诉举报线索多等问题，要深入排查资金收支、合同签订、供应商管理等方面的薄弱环节，及时堵塞管理漏洞，防范类似问题反复发生；建立销号管理制度，跟踪线索后续处理及还款计划落实情况，直至问题解决。

\section{严格考核奖惩，对违规行为严肃追责问责}

各中央企业要全面加强对子企业清欠工作的督促和指导，发挥业绩考核引导作用，推动子企业切实落实工作责任。一是将清欠工作情况和长效机制建设情况纳入企业内部考核，明确集团应付账款和应付票据总体管控目标，针对不同子企业业务类型，分类制定重点子企业应付账款和应付票据管控指标，避免“一刀切”，并严格考核奖惩。二是加强监督检查，将清欠工作纳入内部财务监督检查、内部审计、纪检监察、巡视巡查等工作范围，对中小企业账款支付、拖欠线索核实处理、信息化管控、考核奖惩设置等情况进行检查抽查。三是加强追责问责，对各类拖欠行为始终保持高压态势，对强制中小企业接受非现金结算、应付账款融资确权不及时、拖延验收结算、拖欠投诉线索居高不下、核实处理敷衍了事等行为，要倒查工作落实、长效机制建设、监督检查责任，严肃处理相关责任人员，特别是相关单位主要负责人。国资委将视情况对金额大、反映强烈、影响恶劣的拖欠问题，组织开展调查核实和责任追究工作。四是做好信息公开，严格落实《条例》要求，如实将逾期尚未支付中小企业款项的合同数量、金额等信息纳入国家企业信用信息公示系统企业年度报告，通过系统向社会公示，主动接受社会监督。

各中央企业要扎实做好清欠工作，深入排查拖欠隐患，妥善做好资金安排，及时足额支付民营企业、中小企业账款，确保无分歧欠款“零拖欠”。本通知要求落实的情况要纳入集团清欠工作情况报告，于每季度结束后10日内报送国资委（财务监管与运行评价局）。

~\\

\rightline{国资委}

\rightline{2021年12月17日}